\documentclass[11pt]{article}

\usepackage[a4paper,margin=2.6cm]{geometry}
\usepackage{fontspec}
\setmainfont{Comic Neue}
\setsansfont{Comic Neue}
\usepackage{amsmath,amssymb,amsthm,mathtools}
\usepackage{hyperref}

\newtheorem{theorem}{Theorem}[section]
\newtheorem{proposition}[theorem]{Proposition}
\newtheorem{lemma}[theorem]{Lemma}
\newtheorem{remark}[theorem]{Remark}

\newcommand{\C}{\mathbb C}
\newcommand{\R}{\mathbb R}
\newcommand{\dd}{\,d}
\newcommand{\1}{\mathbf 1}
\newcommand{\calD}{\mathcal D}
\newcommand{\calL}{\mathcal L}
\newcommand{\LambdaU}{\Lambda_U}

\title{The \(h_1\) Boundary Equation from the DtN Strategy}
\author{}
\date{}

\begin{document}
\maketitle

\section{From localization to a real free-boundary problem}

Let \(U\subset\C\simeq\R^2\) be bounded.  In the Bargmann model, the condition that
\(h_1\) be an eigenfunction for the \(h_1\)-window localization operator gives
\begin{equation}\label{eq:h1-localization}
  \int_U (1-\pi |z|^2)e^{-\pi |z|^2}(w-z)e^{\pi w\bar z}\dd A(z)
  =
  \lambda w ,
  \qquad w\in\C .
\end{equation}
Set
\[
  a(r):=(1-\pi r^2)e^{-\pi r^2},
  \qquad
  \rho(r):=(1-4\pi r^2+\pi^2r^4)e^{-\pi r^2}.
\]
Differentiating \eqref{eq:h1-localization} in \(w\) gives
\begin{equation}\label{eq:differentiated-identity}
  \int_U a(|z|)
  \bigl(1-\pi |z|^2+\pi w\bar z\bigr)e^{\pi w\bar z}\dd A(z)
  =
  \lambda .
\end{equation}

Let \(x=(x_1,x_2)\in\R^2\) denote the point \(z=x_1+ix_2\), and define the compactly
supported real distribution
\begin{equation}\label{eq:Omega-def}
  \Omega
  :=
  \rho(|x|)\1_U
  -a(|x|)(x\cdot\nu)\,\delta_{\partial U}
  +\lambda\delta_0 ,
\end{equation}
where \(\nu\) is the outward unit normal.  Equivalently, since
\(\nabla\1_U=-\nu\delta_{\partial U}\),
\[
  \Omega
  =
  \rho(|x|)\1_U+a(|x|)\,x\cdot\nabla\1_U+\lambda\delta_0 .
\]

\begin{lemma}\label{lem:annihilation}
For every \(w\in\C\),
\[
  \bigl\langle \Omega,e^{\pi w\bar z}\bigr\rangle=0.
\]
Since \(\Omega\) is real-valued, also
\[
  \bigl\langle \Omega,e^{\pi \bar w z}\bigr\rangle=0.
\]
\end{lemma}

\begin{proof}
Let \(\phi_w(z)=e^{\pi w\bar z}\).  The radial vector field \(a(|x|)x\) satisfies
\[
  \operatorname{div}(a(|x|)x)
  =
  2(1-3\pi |x|^2+\pi^2|x|^4)e^{-\pi |x|^2},
\]
and therefore
\begin{equation}\label{eq:rho-div-identity}
  \rho(|x|)-\operatorname{div}(a(|x|)x)
  =
  -a(|x|)(1-\pi |x|^2).
\end{equation}
Moreover,
\[
  x\cdot\nabla\phi_w=\pi w\bar z\,\phi_w .
\]
Using the divergence theorem in the distributional form encoded by
\eqref{eq:Omega-def}, we get
\begin{align*}
  \langle\Omega,\phi_w\rangle
  &=
  \int_U \rho(|x|)\phi_w\,\dd A
  -
  \int_{\partial U}a(|x|)(x\cdot\nu)\phi_w\,\dd s
  +\lambda \\
  &=
  \int_U
  \bigl(\rho(|x|)-\operatorname{div}(a(|x|)x)\bigr)\phi_w\,\dd A
  -
  \int_U a(|x|)\,x\cdot\nabla\phi_w\,\dd A
  +\lambda \\
  &=
  -\int_U a(|z|)
  \bigl(1-\pi |z|^2+\pi w\bar z\bigr)e^{\pi w\bar z}\dd A(z)
  +\lambda .
\end{align*}
This vanishes by \eqref{eq:differentiated-identity}.  Taking complex conjugates gives
the second assertion.
\end{proof}

By the same Paley--Wiener/Ehrenpreis--Malgrange divisibility input used in the
Section~4 argument, Lemma~\ref{lem:annihilation} yields a compactly supported real
distribution \(u\) such that
\begin{equation}\label{eq:global-poisson}
  \Delta u=\Omega .
\end{equation}
Thus, away from the origin and in the classical regime where \(u\) is represented by a
function in \(U\) and is extended by zero outside \(U\), \eqref{eq:global-poisson}
means
\begin{equation}\label{eq:BVP}
  \Delta u=\rho(|x|)
  \quad\text{in }U,\qquad
  u=0,\quad
  \partial_\nu u=a(|x|)(x\cdot\nu)
  \quad\text{on }\partial U .
\end{equation}
The point mass \(\lambda\delta_0\) is radial and is irrelevant for the boundary
calculation below as long as \(0\notin\partial U\).

\section{The DtN boundary equation}

Let
\[
  X(x):=(-x_2,x_1),\qquad L:=X\cdot\nabla ,
\]
and set
\[
  v:=Lu .
\]
Since \(L\) commutes with \(\Delta\) and annihilates radial distributions,
\[
  \Delta v=0
  \qquad\text{in }U .
\]
Thus \(v\) is the harmonic function in \(U\) with the boundary trace computed from
\eqref{eq:BVP}.  On \(\partial U\), the zero Dirichlet trace implies
\[
  \nabla u=(\partial_\nu u)\nu=q\nu,
  \qquad
  q:=a(|x|)(x\cdot\nu).
\]
Therefore
\begin{equation}\label{eq:v-boundary-trace}
  v|_{\partial U}
  =
  q(X\cdot\nu).
\end{equation}
If \(\Lambda_U\) denotes the Dirichlet-to-Neumann map for harmonic functions in \(U\),
then
\begin{equation}\label{eq:dtn-side}
  \partial_\nu v
  =
  \Lambda_U\bigl(q(X\cdot\nu)\bigr)
  \qquad\text{on }\partial U .
\end{equation}

\subsection{Parametrized form of the DtN term}

The term in \eqref{eq:dtn-side} can be written explicitly from a boundary
parametrization.  Since \(U\) is simply connected, write
\[
  \partial U=\Gamma,\qquad
  \gamma:\R/L\mathbb Z\to\Gamma
\]
for an arclength parametrization, with outward normal \(\nu(s)\).  Let
\[
  G_0(x,y):=\frac{1}{2\pi}\log|x-y|
\]
be the planar fundamental solution for \(\Delta\).  For a boundary density \(\mu\),
define the double-layer operator
\[
  (D\mu)(x):=
  \int_0^{L}
  \partial_{\nu(t)}G_0(x,\gamma(t))\,\mu(t)\,\dd t .
\]
Its interior trace is
\[
  (D\mu)|_{\partial U}^{\mathrm{int}}
  =
  \Bigl(\frac12 I+K\Bigr)\mu,
\]
where, for \(s\in\R/L\mathbb Z\),
\begin{equation}\label{eq:K-operator}
  (K\mu)(s)
  =
  \operatorname{p.v.}\int_0^{L}
  \frac{(\gamma(t)-\gamma(s))\cdot\nu(t)}
       {2\pi|\gamma(s)-\gamma(t)|^2}
  \,\mu(t)\,\dd t .
\end{equation}
Consequently, for boundary data \(f\), the double-layer density
\(\mu_f\) of the harmonic solution with trace \(f\) is determined by
\begin{equation}\label{eq:mu-equation}
  \Bigl(\frac12 I+K\Bigr)\mu_f=f .
\end{equation}
The normal derivative of the same harmonic function is the hypersingular
operator
\[
  \Lambda_U f=W\mu_f,
\]
where \(W\) may be written in weakly singular form as
\begin{equation}\label{eq:W-operator}
  (W\mu)(s)
  =
  \partial_s
  \operatorname{p.v.}\int_0^{L}
  G_0(\gamma(s),\gamma(t))\,\partial_t\mu(t)\,\dd t .
\end{equation}
Equivalently,
\begin{equation}\label{eq:param-dtn}
  \Lambda_U f
  =
  W\Bigl(\frac12 I+K\Bigr)^{-1}f .
\end{equation}
For the present problem the relevant boundary datum is
\begin{equation}\label{eq:f-def}
  f(s)
  =
  q(s)\,X(\gamma(s))\cdot\nu(s),
  \qquad
  q(s):=
  a(|\gamma(s)|)\,\gamma(s)\cdot\nu(s).
\end{equation}
Thus the DtN term in \eqref{eq:dtn-side} is the explicit nonlocal quantity
\[
  W\Bigl(\frac12 I+K\Bigr)^{-1}
  \bigl(q\,X\cdot\nu\bigr).
\]

We now compute the same normal derivative directly from the boundary jets.  Let
\(T\) be the arclength unit tangent, and use the signed curvature convention
\[
  T'=\kappa\nu .
\]

\begin{lemma}\label{lem:boundary-jets}
On \(\Gamma\),
\[
  D^2u(T,T)=-\kappa q,\qquad
  D^2u(T,\nu)=\partial_s q,\qquad
  D^2u(\nu,\nu)=\rho(|x|)+\kappa q .
\]
\end{lemma}

\begin{proof}
Since \(u=0\) on \(\Gamma\), differentiating twice along arclength gives
\[
  D^2u(T,T)+\nabla u\cdot T'=0.
\]
Using \(T'=\kappa\nu\) and \(\nabla u=q\nu\), this gives
\[
  D^2u(T,T)=-\kappa q.
\]
Tangentially differentiating \(\partial_\nu u=q\) gives
\[
  D^2u(T,\nu)+\nabla u\cdot\nu'=\partial_s q.
\]
Since \(\nu'\) is tangent and \(\nabla u=q\nu\), the second term vanishes.  Finally,
\[
  \rho(|x|)
  =
  \Delta u
  =
  D^2u(T,T)+D^2u(\nu,\nu),
\]
which gives the claimed normal-normal component.
\end{proof}

\begin{proposition}[DtN boundary equation]\label{prop:dtn-equation}
On \(\Gamma=\partial U\),
\begin{equation}\label{eq:dtn-boundary-equation}
  \boxed{
  W\Bigl(\frac12 I+K\Bigr)^{-1}\bigl(q(X\cdot\nu)\bigr)
  =
  (X\cdot T)\,\partial_s q
  +(X\cdot\nu)\bigl(\rho(|x|)+\kappa q\bigr)
  }
\end{equation}
where
\[
  q=a(|x|)(x\cdot\nu).
\]
In arclength coordinates this says the following.  Let \(\mu\) be the solution of
\begin{equation}\label{eq:explicit-mu}
  \frac12\mu(s)
  +
  \operatorname{p.v.}\int_0^L
  \frac{(\gamma(t)-\gamma(s))\cdot\nu(t)}
       {2\pi|\gamma(s)-\gamma(t)|^2}
  \mu(t)\,\dd t
  =
  q(s)\,X(\gamma(s))\cdot\nu(s).
\end{equation}
Then \eqref{eq:dtn-boundary-equation} is exactly
\begin{align}\label{eq:explicit-param-equation}
  \partial_s
  \operatorname{p.v.}\int_0^L
  \frac{1}{2\pi}\log|\gamma(s)-\gamma(t)|\,\partial_t\mu(t)\,\dd t
  &=
  (X(\gamma(s))\cdot T(s))\,\partial_s q(s) \notag\\
  &\quad+
  (X(\gamma(s))\cdot\nu(s))
  \bigl(\rho(|\gamma(s)|)+\kappa(s)q(s)\bigr).
\end{align}
\end{proposition}

\begin{proof}
Directly,
\[
  \partial_\nu v
  =
  \nabla(X\cdot\nabla u)\cdot\nu
  =
  D^2u(\nu,X)+(DX\,\nu)\cdot\nabla u .
\]
The matrix \(DX\) is rotation by \(\pi/2\), so \((DX\,\nu)\cdot\nu=0\).  Since
\(\nabla u=q\nu\), the last term vanishes.  Decomposing
\[
  X=(X\cdot T)T+(X\cdot\nu)\nu
\]
and using Lemma~\ref{lem:boundary-jets}, we obtain
\[
  \partial_\nu v
  =
  (X\cdot T)\,\partial_s q
  +(X\cdot\nu)\bigl(\rho(|x|)+\kappa q\bigr).
\]
Equating this expression with the DtN expression \eqref{eq:dtn-side} proves
\eqref{eq:dtn-boundary-equation}.
\end{proof}

\subsection{The cleaner scalar form}

The equation above looks worse than its actual geometric content.  Write
\[
  \alpha(s):=\gamma(s)\cdot\nu(s),
  \qquad
  \beta(s):=\gamma(s)\cdot T(s),
  \qquad
  r(s):=|\gamma(s)|.
\]
For the standard boundary orientation compatible with the outward normal,
\[
  X(\gamma)\cdot T=\alpha,\qquad X(\gamma)\cdot\nu=-\beta.
\]
Moreover,
\[
  r_s=\frac{\beta}{r},
  \qquad
  \alpha_s=-\kappa\beta.
\]
Since \(q=a(r)\alpha\),
\[
  \partial_s q
  =
  \beta\left(\frac{a'(r)}{r}\alpha-a(r)\kappa\right).
\]
Substituting these identities into
\eqref{eq:dtn-boundary-equation} gives the scalar form
\begin{equation}\label{eq:clean-dtn-equation}
  \boxed{
  \Lambda_U\bigl(-a(r)\alpha\beta\bigr)
  =
  \beta\left(
  \frac{a'(r)}{r}\alpha^2-\rho(r)-2a(r)\kappa\alpha
  \right).
  }
\end{equation}
Thus the layer-potential formula is not the main point; it is only one explicit
realization of the nonlocal operator \(\Lambda_U\).  The boundary equation itself is
the scalar pseudodifferential condition \eqref{eq:clean-dtn-equation}.

Because \(U\) is simply connected, there is also a conformal way to write the same
operator.  If \(\phi:\mathbb D\to U\) is conformal and
\[
  \gamma(\theta):=\phi(e^{i\theta}),
  \qquad
  J(\theta):=|\gamma_\theta(\theta)|,
\]
then, for any boundary datum \(F\),
\begin{equation}\label{eq:conformal-dtn}
  \Lambda_U F(\gamma(\theta))
  =
  \frac{1}{J(\theta)}
  |D_\theta|\bigl(F\circ\gamma\bigr)(\theta)
  =
  \frac{1}{J(\theta)}
  H\partial_\theta\bigl(F\circ\gamma\bigr)(\theta),
\end{equation}
where \(H\) is the periodic Hilbert transform.  Hence
\eqref{eq:clean-dtn-equation} may equivalently be written as
\begin{equation}\label{eq:conformal-clean-dtn}
  \frac{1}{J}|D_\theta|\bigl(-a(r)\alpha\beta\bigr)
  =
  \beta\left(
  \frac{a'(r)}{r}\alpha^2-\rho(r)-2a(r)\kappa\alpha
  \right),
\end{equation}
with all quantities evaluated on \(\gamma(\theta)\).

\begin{remark}[What this does and does not give]
For the Gaussian zero-Neumann problem, applying \(L\) produces a harmonic function
with zero boundary trace, so one gets a pointwise radius equation.  Here the boundary
trace is
\[
  Lu|_{\partial U}=-a(r)\alpha\beta,
\]
which is generally nonzero.  Thus the method does not give an immediate local
rigidity theorem.  What it gives is the closed nonlocal boundary equation
\eqref{eq:clean-dtn-equation}.
\end{remark}

\begin{remark}[Linearized information near centered disks]
The nonlocal equation still has useful content.  Centered disks satisfy it
identically, since \(\beta=0\).  Linearizing around a centered disk of radius \(R\),
with
\[
  r(\theta)=R+\varepsilon h(\theta),
\]
one obtains, for a nonzero Fourier mode \(h_n e^{in\theta}\),
\[
  (1-\pi R^2)(|n|+1-\pi R^2)n\,h_n=0.
\]
Consequently, away from the resonant radii
\[
  \pi R^2=1
  \qquad\text{or}\qquad
  \pi R^2=1+|n|,
\]
the DtN equation is linearly rigid against nonradial perturbations of centered
disks.  At the resonant radii this first-order test leaves the corresponding angular
modes open; higher-order information or an additional identity would be needed there.
\end{remark}

\end{document}
